\chapter{Fazit und Schlussfolgerungen}
\label{chap:fazit}

\section{Ziel und Einordnung des Kapitels}
\label{sec:conclusion_scope}

Kapitel~8 schließt die Arbeit ab, indem die in Kapitel~6 berichteten Ergebnisse zusammengeführt und die in Kapitel~7 diskutierten Einordnungen in prägnante Schlussfolgerungen überführt werden.
Im Unterschied zu Kapitel~6 (deskriptiv) und Kapitel~7 (theoriegeleitet) dient dieses Kapitel der abschließenden Synthese und der expliziten Beantwortung der Forschungsfragen.
Es werden keine neuen Daten eingeführt und keine zusätzlichen Metriken berechnet.

Übergeordnetes Ziel der Arbeit war die Evaluation eines interaktiven KI-Assistenzsystems im Vibe-Coding-Kontext in einer kontrollierten Brownfield-Sandbox.
Im Studiensetup wurde -- konsistent zur Begriffsklärung in Kapitel~1 -- zwischen Assistenzumgebung (Tool), Modell (LLM) und KI-Assistenzsystem als funktionaler Einheit unterschieden.

\section{Beantwortung der Forschungsfragen}
\label{sec:conclusion_rqs}

\subsection[Forschungsfrage 1: Kompetenz und Pipeline-Status]{Forschungsfrage 1: Wie hängt die Entwicklerkompetenz mit dem Erreichen der Pipeline-Status Attempted, \texorpdfstring{Implemented\textsubscript{static}}{Implementedstatic} und Verified zusammen?}
\label{sec:conclusion_rq1}

\textbf{Kurzantwort.}
Attempted und Implemented\textsubscript{static} sind über die Kompetenzgruppen hinweg als Artefaktspuren sichtbar; Verified als testgebundene Stufe wird am seltensten erreicht.

\textbf{Verdichtete Begründung.}
Kapitel~6 trennt die Evidenzebenen Attempted, Implemented\textsubscript{static} und Verified.
Attempted steht für Intention (Chat-Ebene), Implemented\textsubscript{static} für White-Box-Artefakte (Codebasis), Verified für testgebundenes Black-Box-Verhalten und ist damit methodengebunden.
Für die Trichterdarstellung wird in Kapitel~6 zusätzlich Implemented$^\ast$ als abgeleitete Evidenzstufe verwendet.
Die Befunde aus Kapitel~6 und die Einordnung in Kapitel~7 sind damit vereinbar, dass Implementationsfortschritt nicht notwendigerweise in testgebundene Bestätigung übergeht.
Kompetenz dient dabei als Analyseperspektive auf unterschiedliche Integrations- und Validierungsmodi, nicht als Leistungsbewertung.

\subsection[Forschungsfrage 2: Code- und Test-Artefakte]{Forschungsfrage 2: Welche Rolle spielen Code- und Test-Artefakte für die Bewertung der Arbeitsergebnisse eines KI-Assistenzsystems?}
\label{sec:conclusion_rq2}

\textbf{Kurzantwort.}
Code- und Test-Artefakte sind zentral, weil sie die Trennung zwischen sichtbarer Implementierung (Implemented\textsubscript{static}) und testgebundener Bestätigung (Verified) operationalisieren und Chat-Signale methodisch ergänzen.

\textbf{Verdichtete Begründung.}
Die Artefaktanalyse macht Implementationsfortschritt über Quellcodeänderungen, Surrogat-Signale (z.B. Linting) und strukturelle Spuren in der Codebasis nachvollziehbar.
Testgebundene Ergebnisse bilden beobachtbares Verhalten unter den gegebenen Tests ab und sind damit zentral für die Bestätigung im Systemkontext.
Die Pipeline-Logik liefert so eine nachvollziehbare Trennung von Intention, Implementationsartefakt und testgebundener Bestätigung (Kapitel~6) und wird in Kapitel~7 methodisch eingeordnet.
Kompetenz ergänzt diese Sicht als Analyseperspektive, weil sich Artefaktspuren nicht sinnvoll in einem einzelnen, gruppenunabhängigen Signal verdichten lassen.

\subsection[Forschungsfrage 3: Chat-Interaktionsmuster]{Forschungsfrage 3: Welche Chat-Interaktionsmuster treten im Vibe-Coding-Kontext auf, und wie sind sie für die Evaluation einzuordnen?}
\label{sec:conclusion_rq3}

\textbf{Kurzantwort.}
Im Vibe-Coding-Kontext variieren Chat-Interaktionsmuster deutlich; ihre Intensität ist weder als Effizienzmaß noch als Qualitätsmaß geeignet.
Für die Evaluation sind Chat-Muster nur im Zusammenspiel mit Pipeline-Status sowie Code- und Test-Artefakten aussagekräftig.

\textbf{Verdichtete Begründung.}
Die Chat-Artefakte zeigen unterschiedliche Prompt-Strukturen, wechselnde Interaktionsmodi und wiederkehrende Debugging-Schleifen.
Diese Muster sind konsistent damit, dass Interaktion sowohl Exploration als auch Beschleunigung unterstützen kann.
Eine hohe Interaktionsintensität kann dabei gleichermaßen aus Validierung, aus Fehlanpassungen an den Systemkontext oder aus iterativer Integration resultieren.
Entsprechend ist die Einordnung von Chat-Mustern an die Pipeline-Logik zu koppeln.
Attempted macht die Intention sichtbar, Implemented\textsubscript{static} die resultierenden Artefakte und Verified die testgebundene Funktion im System.
Die Kompetenzgruppen liefern hierfür den Kontext als Analyseperspektiven auf Nutzungsformen, ohne dass daraus eine Bewertung einzelner Personen abgeleitet wird.

\section{Zentrale Erkenntnisse der Arbeit}
\label{sec:conclusion_key_findings}

Die Arbeit verdichtet sich in folgenden Befundlinien.

\begin{itemize}
    \item \textbf{Kompetenz als Perspektive auf Nutzungstypen:} Unterschiede zeigen sich in Nutzungsformen des KI-Assistenzsystems, nicht als Wertung von Personen oder Gruppen.
    \item \textbf{Trennung von Intention, Artefakt und Bestätigung:} Attempted, Implemented\textsubscript{static} und Verified bilden unterschiedliche Ebenen desselben Arbeitsprozesses ab und fallen nicht notwendigerweise zusammen.
    \item \textbf{Validierung und Integration als Engpass:} Sichtbare Implementationsartefakte sind häufig erreichbar, während testgebundene Bestätigung die methodisch engste Stufe darstellt.
    \item \textbf{Chat-Interaktion ist kein Ersatz für Ergebnisprüfung:} Interaktionsintensität und Prompt-Formen sind ohne Artefakt- und Testbezug nicht belastbar als Effizienz- oder Qualitätssignale.
    \item \textbf{Grenzen vereinfachter Produktivitätsnarrative:} Produktivität lässt sich in der untersuchten Umgebung nicht auf Generierungsgeschwindigkeit reduzieren, sondern ist entlang von Integrationsaufwand und Validierungsarbeit einzuordnen.
    \item \textbf{Relevanz von Triangulation statt Einzelmetriken:} Belastbare Einordnung entsteht durch das Zusammenspiel von Pipeline-Status, Code- und Test-Artefakten sowie Chat-Mustern.
\end{itemize}

\section{Beitrag der Arbeit}
\label{sec:conclusion_contributions}

\textbf{Wissenschaftlicher Beitrag.}
Die Arbeit liefert eine strukturierte Einordnung der Evaluation eines KI-Assistenzsystems im Vibe-Coding-Kontext, indem Kompetenzgruppen als Analyseperspektiven mit Artefaktspuren und Prozesssignalen zusammengeführt werden.

\textbf{Methodischer Beitrag.}
Die Arbeit operationalisiert Aufgabenerfolg über die Pipeline-Logik (Attempted, Implemented\textsubscript{static}, Verified; sowie Implemented$^\ast$ als abgeleitete Evidenzstufe) und verknüpft diese systematisch mit der Triangulation mehrerer Artefaktquellen.
Damit liegt ein Evaluationsschema vor, das die Trennung zwischen intendierter Bearbeitung, sichtbarer Implementierung und testgebundener Bestätigung explizit und nachvollziehbar macht.

\textbf{Konzeptioneller Beitrag.}
Die Arbeit betont Kompetenz als Perspektive auf Nutzungstypen und verschiebt den Fokus von einer rein toolzentrierten Betrachtung hin zu der Frage, wie KI-Assistenzsysteme in konkreten Entwicklungspraktiken genutzt und geprüft werden.

\textbf{Empirischer Beitrag.}
Die Sandbox-Studie in einer Brownfield-Umgebung liefert eine Artefaktbasis aus Code-, Test- und Chat-Spuren, anhand derer die Pipeline-Status operationalisiert und gruppenbezogen eingeordnet werden können (Kapitel~6/7).

\textbf{Praktischer Beitrag.}
Für Organisationen ergibt sich ein Rahmen zur Einführung und Evaluation von KI-Coding-Tools, der Qualifizierung, Review-Praxis und testbasierte Absicherung als Bestandteile einer kontrollierten Nutzung sichtbar macht.

\section{Abschließendes Fazit}
\label{sec:conclusion_final}

Die Arbeit zeigt, dass die Evaluation KI-gestützter Softwareentwicklung im Vibe-Coding-Kontext eine klare Trennung zwischen Intention, Implementationsartefakt und testgebundener Bestätigung erfordert.
Die Kombination aus Pipeline-Logik, Artefaktanalyse und Chat-Auswertung erlaubt eine konsistente, methodisch abgesicherte Zusammenführung der Befunde (Kapitel~6) und ihrer Einordnung (Kapitel~7).

Im Ergebnis ist der Nutzen im untersuchten Brownfield-Setting nicht als automatische Eigenschaft eines Tools zu lesen, sondern als Zusammenspiel von Integrationsaufwand, Validierungsarbeit und Nutzungstypen im Kompetenzkontext.

\noindent\textbf{Absicherung (Group comparability \& Confounder).} Die Schlussfolgerungen sind an die in Kapitel~6 berichteten Proxys (Pipeline-Status, Code-/Test-Surrogate, Chat-Muster) gebunden; mögliche Konfundierungen und die konservative Lesart werden in Abschnitt~\ref{sec:discussion_scope} (Kapitel~7) begründet. Entsprechend begründen die Befunde keine kausalen Effekte und keine generalisierbaren Rangfolgen zwischen Kompetenzgruppen.
Die Aussagen beziehen sich auf Branches als Analyseeinheiten und auf aggregierte Muster, nicht auf eine Leistungsbewertung einzelner Personen.

\noindent\textbf{Take-Home-Message.} Für KI-gestützte Softwareentwicklung sind singuläre Produktivitäts- oder Proxy-Metriken nur begrenzt belastbar, weil Integrations- und Validierungsarbeit zentral bleibt (Kapitel~6/7). Robust wird die Evaluation durch die konsequente Trennung und Triangulation von Intention (Attempted), Implementationsartefakt (Implemented\textsubscript{static}) und testgebundener Bestätigung (Verified). Kompetenz wird dabei durchgehend als Analyseperspektive genutzt, um unterschiedliche Nutzungs- und Validierungsmodi einzuordnen\,-- ohne Leistungsrangfolgen abzuleiten. Die Kernaussage ist: Belastbare Schlussfolgerungen entstehen erst aus der gemeinsamen Betrachtung dieser drei Ebenen.
